In this chapter, we will simplify the equations describing a system of $N$ particles to the simple case of a binary system. This consists of a body with proper mass $m_1$ and position $\textbf{x}_1$ together with another body of proper mass $m_2$ and position $\textbf{x}_2$. The Lagrangian describing the binary system is obtained by considering $N=2$ in the equation (\ref{eq:NBodySystemLagrangian}), giving 
\begin{align}
L = &\frac{1}{2} m_1 \dot{\textbf{x}}_1^2 + \frac{1}{2} m_2 \dot{\textbf{x}}_2^2 + \frac{G m_1 m_2}{r} + \frac{1}{8} m_1 \frac{(\dot{\textbf{x}}_1^2)^2}{c^2} + \frac{1}{8} m_2 \frac{(\dot{\textbf{x}}_2^2)^2}{c^2} - \frac{1}{2} \frac{G^2}{c^2} \frac{m_1 m _2}{r^2} (m_1 + m_2) \notag \\
& + \frac{1}{2} \frac{G}{c^2} \frac{m_1 m_2}{r} \left[ 3 (\dot{\textbf{x}}_1^2+ \dot{\textbf{x}}_2^2) - 7 \dot{\textbf{x}}_1 \cdot \dot{\textbf{x}}_2 - \frac{(\textbf{x} \cdot \dot{\textbf{x}}_1) (\textbf{x} \cdot \dot{\textbf{x}}_2)}{r^2} \right]  + \Ordof \left( \epsilon^6 \right), \label{eq:PN2BodyLagrangian1}
\end{align}
where we introduce the relative separation vector 
\begin{equation}
\textbf{x} := \textbf{x}_1 - \textbf{x}_2 \label{eq:PN2BodyRelativePosition}
\end{equation}
and its magnitude $r := \left| \textbf{x} \right| = \left| \textbf{x}_1 - \textbf{x}_2 \right|  $.

As usual, we will separate this Lagrangian into two parts, one describing the baricenter and other describing the relative motion. First, we introduce the Tolman mass for each particle as
\begin{align}
\tilde{m}_1 = & m_1 + \frac{1}{2} m_1 \frac{\dot{\textbf{x}}_1^2}{c^2} - \frac{G}{2c^2} \frac{m_1 m_2}{r} \\
\tilde{m}_2 = & m_2 + \frac{1}{2} m_2 \frac{\dot{\textbf{x}}_2^2}{c^2} - \frac{G}{2c^2} \frac{m_1 m_2}{r} 
\end{align}
to define the baricenter position,
\begin{equation}
\textbf{X} = \frac{\tilde{m}_1 \textbf{x}_1 + \tilde{m}_2 \textbf{x}_2}{\tilde{m}_1 + \tilde{m}_2}. \label{eq:PN2BodyBaricenter}
\end{equation}

From these definitions, the acceleration of the baricenter is shown to satisfy
\begin{equation}
(\tilde{m}_1 + \tilde{m}_2) \ddot{\textbf{X}} = \tilde{m}_1 \ddot{\textbf{x}}_1 + \tilde{m}_2 \ddot{\textbf{x}}_2 + \Ordof \left( \epsilon^4 \right).
\end{equation}

Using the Einstein-Infeld-Hoffmann equations to replace the accelerations $\ddot{\textbf{x}}_1$ and $\ddot{\textbf{x}}_2$, it is straightforward to show that the baricenter is not accelerated,
\begin{equation}
\frac{d^2 \textbf{X}}{dt^2} = 0 + \Ordof \left( \epsilon^4 \right),
\end{equation}
and therefore it is possible to describe the dynamics of the binary system with an inertial system attached to the barycenter, in which  $\textbf{X} = 0$. This assumption, together with equations \ref{eq:PN2BodyRelativePosition} and \ref{eq:PN2BodyBaricenter}, let us solve for $\textbf{r}_1$ and $\textbf{r}_2$ in terms of the relative separation,
\begin{align}
\textbf{x}_1 = & \left[ \frac{m_2}{m} + \frac{\mu \Delta m}{2m^2 c^2} \left( \dot{\textbf{x}}^2 - \frac{Gm}{r} \right) \right] \textbf{x} + \Ordof \left( \epsilon^4 \right) \\
\textbf{x}_2 = & \left[ -\frac{m_1}{m}  + \frac{\mu \Delta m}{2m^2 c^2} \left( \dot{\textbf{x}}^2 - \frac{Gm}{r} \right) \right] \textbf{x} + \Ordof \left( \epsilon^4 \right),
\end{align}
where we introduce $m=m_1 + m_2$, $\mu = {m_1 + m_2}{m}$ and $\Delta m= m_1 - m_2 $. The corresponding velocities satisfy the relations
\begin{align}
\dot{\textbf{x}}_1 = & \frac{m_2}{m} \dot{\textbf{x}} + \Ordof \left( \epsilon^2 \right) \\
\dot{\textbf{x}}_2 = & - \frac{m_1}{m} \dot{\textbf{x}} + \Ordof \left( \epsilon^2 \right) .
\end{align}

In the baricenter frame, the Lagrangian (\ref{eq:PN2BodyLagrangian1}) describes the relative motion of the particles in the binary system, and takes the form
\begin{align}
L = &\frac{1}{2} \mu \dot{\textbf{x}}^2 + \frac{G \mu m}{r} + \frac{1}{8} \mu \left( 1 - \frac{3\mu}{m} \right) \frac{(\dot{\textbf{x}}^2)^2}{c^2} - \frac{1}{2} \frac{G^2}{c^2} \frac{\mu m^2}{r^2} \notag \\
& + \frac{1}{2} \frac{G}{c^2} \frac{\mu m}{r} \left[ \left( 3 + \frac{\mu }{m}\right) \dot{\textbf{x}}^2 + \frac{\mu }{m} \left( \frac{\textbf{x} \cdot \dot{\textbf{x}}}{r} \right)^2 \right]  + \Ordof \left( \epsilon^6 \right) .
\end{align}

Note that the first two terms in this Lagrangian correspond to the Newtonian contribution while the remaining terms are post-Newtonian corrections. The Euler-Lagrange equations associated to this Lagrangian are
\begin{align}
\ddot{x} ^k = &-\frac{Gm}{r^3} x^k + \frac{Gm}{c^2 r^3} \left[  \frac{2Gm}{r} \left( 2 + \frac{\mu }{m} \right) - \left( 1 + \frac{3\mu }{m} \right) (\dot{\textbf{x}}^2)^2 + \frac{3}{2} \frac{\mu }{m} \left( \frac{\textbf{x} \cdot \dot{\textbf{x}}}{r} \right)^2   \right] x^k \notag \\
& + \frac{2Gm}{c^2 r^3} (\textbf{x} \cdot \dot{\textbf{x}}) \left( 2 - \frac{\mu }{m} \right) \dot{x}^k  + \Ordof \left( \epsilon^6 \right), 
\end{align}
or, in vector notation, 
\begin{align}
\ddot{\textbf{x}} = &-\frac{Gm}{r^3} \textbf{x} + \frac{Gm}{c^2 r^3} \left[  \frac{2Gm}{r} \left( 2 + \frac{\mu }{m} \right) - \left( 1 + \frac{3\mu }{m} \right) (\dot{\textbf{x}}^2)^2 + \frac{3}{2} \frac{\mu }{m} \left( \frac{\textbf{x} \cdot \dot{\textbf{x}}}{r} \right)^2   \right] \textbf{x} \notag \\
& + \frac{2Gm}{c^2 r^3} (\textbf{x} \cdot \dot{\textbf{x}}) \left( 2 - \frac{\mu }{m} \right) \dot{\textbf{x}} + \Ordof \left( \epsilon^6 \right). \label{eq:PNBinaryRelativeAcc}
\end{align}

\section{Post-Newtonian Motion of the Binary System}

The zeroth order motion of the binary system corresponds to the Keplerian orbit. Choosing coordinates such that the orbit  lies in the $z=0$ plane, the Keplerian conic can be written as
\begin{equation}
\textbf{x} = r \cos f \textbf{e}_x + r \sin f \textbf{e}_y , \label{eq:PNBinaryRelativePosition}
\end{equation}
where
\begin{equation}
r = \frac{p}{1+e \cos f} \label{eq:PNBinaryRadius}
\end{equation}
where $p$ is the semi-latus rectum and $f=\varphi - \omega_0$ is the true anomaly. Now we set, without losing generality, the initial argument of the pericenter as $\omega_0 =0$. The zeroth order solution has a conserved angular momentum that gives the relation
\begin{equation}
r^2 \frac{d\varphi}{dt} = \sqrt{Gmp}.
\end{equation}

In order to introduce the post-Newtonian contributions to the motion of the binary system, we will introduce a small correction term, $\delta \ell \sim \epsilon \ll 1$, to the angular momentum equation
\begin{equation}
r^2 \frac{d\varphi}{dt} = \sqrt{Gmp} (1 + \delta \ell), \label{eq:PNBinaryCorrectedAngMom}
\end{equation}
and a small correction, $\frac{\delta \dot{\textbf{x}}}{c} \sim \epsilon \ll 1$, to the velocity
\begin{equation}
\dot{\textbf{x}} = \frac{d \textbf{x}}{dt}  = \sqrt{\frac{Gm}{p}} \left[ -\sin \varphi \textbf{e}_x + (e+\cos \varphi)\textbf{e}_y \right] + \delta \dot{\textbf{x}}. \label{eq:PNBinaryCorrectedVel}
\end{equation}

Replacing (\ref{eq:PNBinaryCorrectedAngMom}) into the angular momentum definition $r^2 \frac{d\varphi}{dt}  = \left| \textbf{x} \times \dot{\textbf{x}}\right|$, we obtain a differential equation for the correction $\delta \ell$,
\begin{equation}
\frac{d}{dt} \left( \delta \ell \right) = \frac{1}{\sqrt{Gmp}} \left| \textbf{x} \times \ddot{\textbf{x}}\right|.
\end{equation}

Using equations (\ref{eq:PNBinaryRelativeAcc}) for the relative acceleration and (\ref{eq:PNBinaryCorrectedVel}) for the velocity, we obtain 
\begin{equation}
\frac{d}{dt} \left( \delta \ell \right) = - \frac{2Gm}{c^2 p} e \left( 2 - \frac{\mu}{m} \right) \frac{d \cos \varphi}{dt},
\end{equation}
which gives immediately the correction to the angular momentum as
 \begin{equation}
 \delta \ell  = - \frac{2Gm}{c^2 p} e \left( 2 - \frac{\mu}{m} \right) \cos \varphi . \label{eq:PNBinaryCorrectiontoAngMom}
\end{equation}

Replacing these results into equation (\ref{eq:PNBinaryRelativeAcc}) and integrating, we find the relative velocity in binary system as
\begin{equation}
 \dot{\textbf{x}} = v_x \textbf{e}_x + v_y \textbf{e}_y
 \end{equation} 
 with
 \begin{subequations}\label{eq:PNBinaryVelocityComponents}
 \begin{align}
 v_x = &-\sqrt{\frac{Gm}{p}} \sin \varphi
 - \frac{1}{c^2} \left( \frac{Gm}{p} \right)^{3/2}  \left[ 3e\varphi - \left( 3 - \frac{\mu}{m} \right)\sin \varphi + \left( 1 + \frac{21\mu}{8m} \right) e^2 \sin \varphi \right. \notag \\  
 & \left. -\frac{1}{2} \left( 1 - \frac{2\mu}{m} \right)e \sin 2 \varphi + \frac{\mu}{8m} e^2 \sin 3 \varphi\right] + \Ordofee \\
 v_y = &-\sqrt{\frac{Gm}{p}} \left( e + \cos \varphi \right ) - \frac{1}{c^2} \left( \frac{Gm}{p} \right)^{3/2} \left[  \left( 3 - \frac{\mu}{m} \right)\cos \varphi + \left( 3 - \frac{31\mu}{8m} \right) e^2 \cos \varphi  \right .  \notag \\
 & \left. + \frac{1}{2} \left( 1 - \frac{2\mu}{m} \right)e \cos 2 \varphi - \frac{\mu}{8m} e^2 \cos 3 \varphi\right] + \Ordofee
 \end{align}
 \end{subequations}
 
\subsection{Pericenter Shift}

\begin{exercise}
\textbf{Post-Newtonian Correction to the Orbit of the Binary System}

Use equations (\ref{eq:PNBinaryRelativePosition}),(\ref{eq:PNBinaryCorrectedAngMom}), (\ref{eq:PNBinaryCorrectiontoAngMom}) and (\ref{eq:PNBinaryVelocityComponents}), together with the relation
\begin{equation}
\frac{d}{d\varphi} \left( \frac{1}{r} \right) = - \left( r^2 \frac{d\varphi}{dt} \right)^{-1} \frac{\textbf{x} \cdot \dot{\textbf{x}}}{r},
\end{equation}
to obtain the post-Newtonian corrected orbit of the binary system,
\begin{align}
r = p & \left[ 1 + e \cos \varphi + \frac{Gm}{c^2 p}  \left[  -\frac{1}{2} \left( 7 - \frac{2\mu}{m} \right) e  + \frac{\mu}{4m} e^2 + 3 e \varphi \sin \varphi \right. \right. \notag \\
& \left. \left. + \frac{1}{2} \left(7 - \frac{2\mu}{m} \right) e \cos \varphi - \frac{\mu}{4m} e^2 \cos 2\varphi   \right] + \Ordofee \right]^{-1} , \label{eq:PNBinaryCorrectedConic}
\end{align}
which replaces the conic equation (\ref{eq:PNBinaryRadius}).
\end{exercise}

In order to obtain an estimate for the motion of the pericenter of a binary system, it is necessary to analyze the post-Newtonian contributions in equation (\ref{eq:PNBinaryCorrectedConic}) which are very similar to those obtained for the motion of a particle in Schwarzschild's background in equation (\ref{eq:FirstOrderApproxGeneralOrbit}). Just as studied in Section \ref{sec:PerihelionSchwarzschild}, it is clear that the secular contribution to the pericenter motion is given by the third correction term, $ 3 \frac{Gm}{c^2 p} e \varphi \sin \varphi $, and it gives a shift of
\begin{equation}
\Delta \omega =  \frac{6 \pi Gm}{c^2 p} \label{eq:BinarySystemPerihelionAdvance}
\end{equation} 
per orbit. This result looks exactly as Schwarzschild's shift, with just one difference: the mass $M$, appearing in equation (\ref{eq:SchwarzschildPerihelionAdvance}), corresponds to the central object while the mass $m$, appearing in (\ref{eq:BinarySystemPerihelionAdvance}), is the sum of the masses of the two bodies in the binary system.
